\par\bigskip
\begin{center}\large\textbf{VI. Die Jahre 1801 bis 1810.}\end{center}
\begin{center}\textls[100]{a) Bis zur Berufung nach Göttingen 1807. Allgemeines.}\\
\textls[100]{Zahlentheoretische Anwendungen der elliptischen Funktionen.}\end{center}
\noindent\hspace*{\parindent}Im Vergleich zu den sich auf erstaunlich kurze Zeiträume zusammendrängenden
Resultaten, über die wir in den vorhergehenden Abschnitten zu
berichten hatten, erscheint der Ertrag der analytischen Forschungen, die \textsc{Gauss}
in den folgenden Jahren angestellt hat, nicht allzu reichlich. Wir werden
sehen, dass er namentlich zu Anwendungen seiner neuen Transzendenten auf
Probleme der Zahlentheorie übergegangen ist; inwieweit freilich in diese Jahre
noch Vorarbeiten für später veröffentlichte oder wenigstens aufgezeichnete
Untersuchungen fallen, lässt sich nicht mit Sicherheit feststellen, zumal auch
die Aufzeichnungen des \textit{Tagebuchs} immer spärlicher werden und bald ganz
versiegen. Der Grund hierfür geht auf ein Ereignis zurück, das auf \textsc{Gauss’}
wissenschaftliche und persönliche Entwicklung von ausschlaggebender Bedeutung
wurde, nämlich auf die Entdeckung des ersten kleinen Planeten, der
Ceres, in der Neujahrsnacht von 1800 auf 1801\begingroup\renewcommand{\thefootnote}{165a)}\footnote{Es mag hier darauf hingewiesen werden, dass \textsc{Gauss} die Wende des \textls[100]{Jahrhunderts} nicht für diese Nacht angesetzt zu haben scheint; er schreibt an \textsc{Bolyai} am 16. Dezember 1799 (\textit{Briefwechsel} S. 37) in bezug auf die bevorstehende Jahreswende: »der letzte Dezember, der wenigstens der letzte Tag sein wird, wo wir siebzehnhundert nennen, wenn gleich mikrologische Ausleger das Ende des Jahrhunderts noch ein Jahr weiter hinaussetzen .\ .\ .«.}\addtocounter{footnote}{-1}\endgroup\ durch \textsc{Josef Piazzi} in Pa-\par
\clearpage\SourceLeaf{203}\ArithmeticHead{118}{L}{L. SCHLESINGER, ÜBER GAUSS’ FUNKTIONENTHEORETISCHE ARBEITEN.}
\noindent lermo\footnote{Näheres siehe bei \textsc{Brendel}, Werke XI, 2, Abhandlung 3, S. 150—155; vergl. auch die Bemerkungen \textsc{Brendels} zu den Tagebuchnotizen 119.—122., Werke X, 1, S. 561—563.}. Die Aufgabe der Wiederauffindung dieses Himmelskörpers interessierte
\textsc{Gauss} auf das lebhafteste, er entwickelte bald eigentümliche Methoden
für die Bahnbestimmung (siehe die Tagebuchnotizen 119., 121. vom September,
Oktober 1801), auf Grund deren kurz hintereinander v. \textsc{Zach} und \textsc{Olbers} die
Ceres wiederfanden. Dies trug mit einem Male \textsc{Gauss’} Namen auch in weitere
Kreise; schon am 31. Januar 1802 wählte die Petersburger Akademie den
jungen Forscher zum Korrespondenten\begingroup\renewcommand{\thefootnote}{166a)}\footnote{Der Sekretär \textsc{Fuss} schlug \textsc{Gauss} vor, als »vorteilhaft bekannt durch seine analytischen Arbeiten und neuerdings durch seine Rechnungen über den neuen Planeten Ceres, denen man es verdankt, diesen wieder aufgefunden zu haben«, siehe W. \textsc{Stieda}, Jahrbücher für Kultur und Geschichte der Slaven, N. F. III, 1927, S. 79 ff. besonders S. 90.}\addtocounter{footnote}{-1}\endgroup.

In dem mit der Notiz 121. (Oktober 1801) unterbrochenen \textit{Tagebuch}
schreibt \textsc{Gauss} (Notiz 122.): »In den folgenden Jahren 1802, 1803, 1804 haben
astronomische Beschäftigungen den grössten Teil meiner Musse in Anspruch
genommen, insbesondere die Rechnungen über die Theorie der neuen Planeten.
So kam es, dass in diesen Jahren der vorliegende Katalog vernachlässigt
wurde. — Die Tage, an denen es mir gegeben war, irgend etwas zur Förderung
der reinen Mathematik beitragen zu können, sind meinem Gedächtnis
entschwunden.« — Neben diesen sind noch andere, mehr persönliche Verhältnisse
der stillen Weiterarbeit an den analytischen Problemen hindernd entgegengetreten.
Wir sind in der glücklichen Lage, darüber \textsc{Gauss’} eigenen
Bericht wiedergeben zu können, in der Form von Briefauszügen an seinen
vertrauten Jugendfreund \textsc{Bolyai}. »Ich lebe«, schreibt er am 3. Dezember 1802,
(\textit{Briefwechsel} S. 45) »seitdem ich 1800 Ostern Helmstedt wieder verlassen habe,
bisher beständig in Braunschweig hauptsächlich für meine Göttinnen, die Wissenschaften.
Bis im Sommer 1801 hat mich die Arbeit an meinem grossen Werke
beschäftigt, welches Michael 1801 herausgekommen ist. Es führt den Titel:
\textit{Disquisitiones arithmeticae}, Lipsiae in Commissis apud G. \textsc{Fleischer} junior,
1801 .\ .\ . Seitdem haben mich besonders die beiden neuen Planeten Ceres
und Pallas beschäftigt. .\ .\ .\ . Liesest Du v. \textsc{Zachs} Monatliche Correspondenz,
so wirst Du schon daraus wissen (Maistück 1802), dass meine Arbeit über die
Ceres die Ursache einer ansehnlichen Verbesserung meiner äusseren Lage ist:
unser grossmütiger Fürst hat mich durch eine Pension von 400 Rthl. vor der\par
\clearpage\SourceLeaf{204}\ArithmeticHead{119}{R}{VIa. VON 1901 BIS ZUR BERUFUNG NACH GÖTTINGEN.}
\noindent Hand in eine unabhängige, sorgenfreie Lage gesetzt. Jedoch ist es noch sehr
ungewiss, ob ich diese Lage noch lange geniessen werde. Gerade jetzt stehe
ich in Unterhandlung wegen eines Rufes nach St. Petersburg als Direktor der
Kais. Sternwarte und ausserdem ist sogar eine Aussicht da, dass ich vielleicht
zwischen diesem Rufe und einem andern die Wahl haben könnte. .\ .\ .\ . Ich
stehe überhaupt mit Göttingen in weniger Verbindung, nur erst vor ein Paar
Wochen hat mich die Sozietät der Wissenschaften zu ihrem Korrespondenten
ernannt.« Mit der Andeutung über den andern Ruf, zwischen dem und Petersburg
er die Wahl haben könnte, ist die Möglichkeit gemeint, die sich
für \textsc{Gauss} eröffnete, an die neuzuerbauende Göttinger Sternwarte berufen zu
werden, wofür sich \textsc{Olbers} — mit dem \textsc{Gauss} seit Januar 1802 im Briefwechsel
stand — bei dem Göttinger Professor \textsc{Heeren} sehr lebhaft eingesetzt
hatte. — Weiter schreibt er an \textsc{Bolyai} am 20. Juni 1803 (\textit{Briefwechsel} S. 55):
»Die Vokation nach St. Petersburg hat mich nicht von hier weggezogen, unser
Herzog liess mich nicht fort und hat mir meine hiesige Lage noch angenehmer
gemacht. Ich habe sogar Hoffnung zu einer kleinen hiesigen Sternwarte
— falls der leidige Krieg nicht von neuem unsere Projekte hemmt —.
Astronomie und Reine Grössenlehre sind einmal die magnetischen Pole, nach
denen sich mein Geisteskompass immer wendet. .\ .\ .\ . Ich reise in diesen
Tagen nach Bremen zu einem Besuch bei Dr. \textsc{Olbers}, dessen Freundschaft
ich, ohne ihn bisher persönlich zu kennen, durch Briefe kultiviert habe.« Und
am 28. Juni 1804 (\textit{Briefwechsel} S. 61): »Seit meinem letzten Briefe an Dich
habe ich eine Menge neuer zum Teil äusserst interessanter Bekanntschaften gemacht.
Einen Monat bin ich in Bremen und Lilienthal bei \textsc{Olbers}, einem der
allerliebenswürdigsten Menschen, die ich kenne, und \textsc{Schröter}, und vier Monat
bin ich in Gotha gewesen. Auch hier hat sich der Kreis meiner Bekannt-
und Freundschaften sehr erweitert. Die schönste aber ist die eines herrlichen
Mädchens, ganz so, wie ich mir immer eine Gefährtin meines Lebens gewünscht
habe. .\ .\ .\ .\ .\ . Dass das Nationalinstitut in Paris und die königl.
Sozietät der Wissenschaften mir unlängst die Ehre erzeigt haben, mich unter
ihre auswärtigen Mitglieder aufzunehmen, weisst Du vielleicht aus öffentlichen
Nachrichten. Erhält mir die Vorsehung ihre mir bisher verliehenen Güter
und schenkt mir dazu noch das höchste häusliche Glück, so sehne ich mich
noch nicht aus diesem Erdenleben nach dem Himmel.« Am 25. November\par
\clearpage\SourceLeaf{205}\ArithmeticHead{120}{L}{L. SCHLESINGER, ÜBER GAUSS’ FUNKTIONENTHEORETISCHE ARBEITEN.}
\noindent 1804 heisst es weiter (\textit{Briefwechsel} S. 79): »Deinen Brief vom 16. September
habe ich zwar schon in der Mitte Oktobers erhalten, indessen bin ich nicht
ohne Entschuldigung, dass ich ihn erst in diesem Monat beantworte. Erstens
wirst Du wahrscheinlich aus öffentlichen Nachrichten längst wissen, dass Anfang
September d. J. in Lilienthal von \textsc{Harding}, einem sehr lieben persönlichen
Freunde von mir, ein dritter neuer Planet entdeckt ist, der den Namen
Juno erhalten hat, und der mir umso überhäuftere Arbeit macht, da ich nicht
nur, wie bei Ceres und Pallas alle Rechnungen darüber allein über mich
nehme, sondern ihn auch selbst beobachte, so oft das Wetter es nur erlaubt,
da ich jetzt sowohl mit Instrumenten versehen bin, als auch bereits hinlängliche
praktische Fertigkeit im Observieren erworben habe. .\ .\ .\ .\ . Lege Dich
doch auch etwas auf die praktische Astronomie; sie ist nach meinem Gefühle,
nächst den Freuden des Herzens und der Beschauung der Wahrheit in der
reinen Mathesis, der süsseste Genuss, den wir auf Erden haben können. Ein
zweiter noch wichtigerer Grund aber, warum ich meine Antwort etwas verschob,
war, weil ich sie nicht gern eher geben wollte, bis ich Dir etwas von
mir erzählen könnte. Das kann ich jetzt, liebster \textsc{Bolyai}. Seit drei Tagen
ist der für diese Erde fast zu himmlische Engel meine Braut. Ich bin überschwenglich
glücklich.« Dann schreibt er am 20. Mai 1808 schon aus Göttingen:
»Meine Lage in Braunschweig hatte ich von jeher nur als eine interimistische
betrachtet, die sich über kurz oder lang verändern müsste. Dass
aber solche Katastrophen mich von da so bald wegtreiben würden, ahndete
ich freilich nicht. Du kennst die unglückliche Geschichte des Herbstes 1806.
Wenige Tage vorher noch im Genuss aller Segnungen des Friedens, sahen
wir auf einmal unsere Fluren zum Schauplatz des Krieges werden, sahen wir
unseren geliebten Fürsten tödtlich verwundet, kaum ein Paar Tage Ruhe in
seinem Lande findend, den Verfolgungen der Feinde fliehend, um bald in
fremder Erde eine Ruhestatt zu finden\footnote{Vergl. oben S. 63, Fussnote \textsuperscript{82)}.}. .\ .\ .\ .\ . Tausende verloren in diesem
Wechsel der Dinge sogleich alles; mir persönlich wurden aber die Folgen
nicht gleich sehr fühlbar. Auch geschahen mir von mehreren Orten her Anträge,
die mich über die Besorgnisse wegsetzten. Ich nahm den Ruf nach
Göttingen als Professor Astronomiae und erster Direktor der Sternwarte an
(mein Kollege ist \textsc{Harding} .\ .\ .\ .) und ich bin seit einem halben Jahre hier.\par
\clearpage\SourceLeaf{206}\ArithmeticHead{121}{R}{VIa. VON 1901 BIS ZUR BERUFUNG NACH GÖTTINGEN.}
\noindent Ich hoffe hier ganz glücklich zu leben. Der König\begingroup\renewcommand{\thefootnote}{167a)}\footnote{\textsc{Jérôme Bonaparte}, 1807—1813 König von Westfalen.}\addtocounter{footnote}{-1}\endgroup\ hat Hoffnung gegeben,
dass der Bau der neuen Sternwarte, der schon 1803 anfing, fortgesetzt werden
soll, dann bleibt mir wenig zu wünschen übrig. Auch die äussern Verhältnisse
sind ganz gut; ich bin mit 1000 Rt. angestellt, auf Nebenverdienst
durch Kollegia kann ich freilich wenig rechnen, zumal jetzt, wo sozusagen
ganz Europa zur Bettlerin geworden ist. In meinem Hause lebe ich sehr
glücklich. Seit dem 9. Oktober 1805 bin ich verheiratet .\ .\ .\ ., seit 21. August
1806 Vater eines allerliebsten Jungen \textsc{Josef} und seit 29. Februar 1808 Vater
eines Mädchens \textsc{Wilhelmine} .\ .\ .\ .\ .\ . Über die Planeten wirst Du vieles in
einem grossen Werke von mir finden, das schon gedruckt wird: \textit{Theoria motus
corporum coelestium circa solem revolventium}«\footnote{Im Druck heisst es statt \textit{circa solem revolventium}, \textit{in sectionibus conicis solem ambientium}; das Werk erschien 1809 bei \textsc{Friedrich Perthes} und I. H. \textsc{Besser} in Hamburg; siehe Werke VII.}. Und endlich am 2. September
1808: »Meine wissenschaftlichen Beschäftigungen gehen ihren Gang fort; der
Druck meines astronomischen Werks, in welchem alles, was sich auf ihre
Bewegung bezieht, vollständig abgehandelt wird, ist über die Hälfte vollendet.
Meine fortgesetzten Untersuchungen aus der höhern Arithmetik denke ich
vorerst grösstenteils als einzelne, in die Commentationes unsrer Sozietät einzurückende
Abhandlungen auszuarbeiten, die vielleicht in Zukunft einmal zu
einem zweiten Bande meiner \textit{Disquis. Arithm.} gesammelt werden können; zwei
davon sind bereits der Sozietät übergeben und eine schon abgedruckt im
XVI. Bande der Göttingischen Commentationes. Von Personen, die jenes
Werk mit Erfolg studiert hätten, kenne ich bis jetzt nur wenige; oben an
steht die Demoiselle \textsc{Sophie Germain} in Paris .\ .\ .\ ., \textsc{Lagrange} hat einen Paragraph
in der neuen Auflage seines \textit{Traité de la résolution numérique des équations}
kommentiert, indess auf eine Art, der ich nicht in allen Stücken meinen
Beifall geben kann. Ein vieljähriger Freund von mir, und mein erster Lehrer
in der Mathematik, \textsc{Bartels}, der .\ .\ .\ .\ . jetzt Professor der Mathematik in
Kasan in Russland geworden ist, hat .\ .\ . jenes Werk auch fleissig studiert
und ist in die meisten Materien eingedrungen. .\ .\ .\ .\ .\ .\ . Merkwürdig ist es
immer, dass alle diejenigen, die diese Wissenschaft ernstlich studieren, eine
Art Leidenschaft dafür fassen. .\ .\ .\ . Wahrlich, es ist nicht das Wissen, sondern
das Lernen, nicht das Besitzen, sondern das Erwerben, nicht das Da-Sein,\par
\HistoricalFoot{X2 Abh. 2.}{16}
\clearpage\SourceLeaf{207}\ArithmeticHead{122}{L}{L. SCHLESINGER, ÜBER GAUSS’ FUNKTIONENTHEORETISCHE ARBEITEN.}
\noindent sondern das Hinkommen, was den grössten Genuss gewährt. Wenn ich eine
Sache ganz ins Klare gebracht und erschöpft habe, so wende ich mich davon
weg, um wieder ins Dunkle zu gehn; so sonderbar ist der nimmersatte Mensch,
hat er ein Gebäude vollendet, so ist es nicht um nun ruhig darin zu wohnen,
sondern um ein anderes anzufangen. So, stelle ich mir vor, muss einem Welteroberer
zu Mute sein, der nachdem ein Königreich kaum bezwungen ist,
schon wieder nach andern seine Arme ausstreckt«\begingroup\renewcommand{\thefootnote}{168a)}\footnote{Mit diesem Briefe gerät der Briefwechsel zwischen \textsc{Gauss} und \textsc{Bolyai} für viele Jahre ins Stocken. \textsc{Gauss} nächster Brief datiert vom 6. März 1832 (siehe Werke VIII, S. 220), dann hat \textsc{Gauss} nur noch zweimal, am 23. Oktober 1836 und am 20. April 1848 an den Jugendfreund geschrieben.}\addtocounter{footnote}{-1}\endgroup.
